\documentclass[a4paper]{ctexart}
\usepackage[margin=1in]{geometry}

\title{第一讲课后练习}
\author{PB24000150 李欣宸}
\date{\today}

\begin{document}

\maketitle

以录取成功率预测为例：

\begin{itemize}
    \item 机器学习可以用来预测申请学校录取的成功率；在科大，「飞跃」是一个源远流长的话题，一个良好的「飞跃」成败预测器可以指导学生更好地选择申请的课题组和润色简历；
    \item 对于这个问题，GPA，本科专业，对方学科方向，专业排名，对方学校排名，科研经历，竞赛经历，是否少院/年龄，性别等都可以作为特征；此外，科大申请同一课题组的其他同学成功/失败情况也可以作为特征；
    \item 对于这个问题上深度学习模型相比线性回归模型所具有的优势，我的想法有：
    \begin{itemize}
        \item 线性回归模型只能建模线性关系，而深度学习模型可以建模非线性关系；申请学校是否录取这个问题中，可能存在一些复杂的非线性关系，比如对方学科方向和本科专业之间的契合程度可能会影响录取成功率，这种关系可能很难用线性回归模型来捕捉； 深度学习模型可以建模更加复杂的关系（比如对方学科方向和本科专业之间的契合程度）而申请学校是否录取并不是严格的线性决策边界（虽然仅仅就这个问题而言，决策边界其实是很接近线性的，可能很适合线性回归模型，但是因为这题要求回答深度学习模型的好处，所以有这个理由）；
        \item 深度学习可以利用线性回归模型难以利用的特征，比如，推荐信的全文很难变成线性回归可用的特征（虽然可以把大量推荐信的嵌入向量进行 PCA 取前三主成分也不失为一种方法）；
        \item 最重要的，深度学习方法应用在问题中，可以更方便地去利用其他来源的一些先验知识；这些问题实际上可以使用一些深度学习与集成学习相结合的策略，比如将对方学校、专业这种难以直接应用线性回归的特征利用深度学习的文本嵌入模型转化为低维度的特征向量；然后再用集成学习（如 XGBoost）等策略建立模型；这样实际利用了深度学习在文本语料库中训练得到的关于大学层次的大量数据。
    \end{itemize}
\end{itemize}

如果刚才这个问题偏题了的话，一个不太可能偏题的组合：

\begin{itemize}
    \item 校园门禁人脸识别；
    \item 人脸图像，面部各个关键点位置；
    \item 决策边界复杂，深度学习的良好表征能力可以良好刻画这种复杂关系；深度学习模型可以端到端，代码简洁；深度学习可以利用其他图像中学到的先验知识，进行迁移学习，比如把图像基座模型改造成人脸识别模型；
\end{itemize}

\end{document}